\chapter{\MakeUppercase{Technical Specification}}
\justifying
\section{Requirements}

\subsection{Functional Requirements}
\paragraph{} The below functional requirements define the fundamental capabilities of the adaptive gameplay system:

\begin{enumerate}
	\item \textbf{FR-01: Game Mechanics:} The system shall provide a fully playable space shooter game with player movement (WASD/Arrow keys), shooting (Spacebar), enemy ships, and asteroids.
	\item \textbf{FR-02: Real-Time Monitoring:} The system shall continuously monitor player performance metrics such as accuracy ratio, kill rate, player health, damage rate, and length of survival time in game.
	\item \textbf{FR-03: Dynamic Difficulty Adjustment:} The system shall dynamically adjust game difficulty settings (enemy speed, enemy health, spawn delay, attack intensity) in real time according to the inference of the \ac{ppo} model.
	\item \textbf{FR-04: Model Training:} The system shall provide a training pipeline for the \ac{ppo} agent using a simulated game environment.
	\item \textbf{FR-05: Data Logging:} The system shall log gameplay data (player metrics, difficulty parameters, timestamps) to a FastAPI backend.
	\item \textbf{FR-06: Analytics Dashboard:} The system shall display real-time visualizations of player performance and difficulty trends through an interactive React dashboard.
	\item \textbf{FR-07: Model Persistence:} The system shall save and load trained model weights (\texttt{dda\_ppo\_3d.pth}) during live gameplay for inference.
\end{enumerate}

\subsection{Non-Functional Requirements}
\begin{enumerate}
	\item \textbf{NFR-01: Frame Rate:} The game shall maintain a stable frame rate (minimum 30 FPS) during gameplay with the \ac{dda} inference enabled.
	\item \textbf{NFR-02: Inference Latency:} The \ac{ppo} model inference shall complete within 50 milliseconds to prevent significant lag during gameplay.
	\item \textbf{NFR-03: Training Stability:} The training pipeline shall produce convergent policies without catastrophic divergence.
	\item \textbf{NFR-04: Modularity:} The system architecture shall support independent modification of game logic, \ac{rl} components, backend, and frontend.
	\item \textbf{NFR-05: Portability:} The system shall run on any platform supporting Python 3.8+ and modern web browsers.
	\item \textbf{NFR-06: Usability:} The game controls shall be intuitive and responsive, requiring no previous experience with computer games.
\end{enumerate}

\section{Architecture Rationale}
\paragraph{} The selection of the technology stack for this project was driven by a need for high-performance execution, seamless integration of machine learning models, and robust real-time data handling. This section details the fundamental rationale behind the architectural choices and the complexity of the data pipeline.

\subsection{Technology Stack Justification}
\paragraph{} \textbf{Game Engine (Pygame):} While industry-standard engines like Unity or Unreal Engine offer powerful built-in physics and rendering capabilities, Pygame was selected for this project due to its bare-metal exposure to game loop mechanics. In Pygame, developers have direct, unfiltered access to the event queue, frame rendering delta time, and memory state. This low-level access is critical for reinforcement learning, as it allows the environment wrapper (via Gymnasium) to predictably freeze the game state, inject agent actions, and extract high-precision arrays of object coordinates without the asynchronous overhead typical of heavier commercial game engines.
\paragraph{} \textbf{Machine Learning Framework (PyTorch):} PyTorch was selected over TensorFlow due to its dynamic computational graph architecture. \ac{dda} via \ac{ppo} requires rapid prototyping of network shapes, particularly when alternating between standard Multi-Layer Perceptrons (MLPs) and Long Short-Term Memory (LSTM) recurrent networks. Furthermore, PyTorch natively supports seamless tensor bridging into NumPy arrays, minimizing the serialization latency when passing game state floats into the neural network during live gameplay.
\paragraph{} \textbf{Backend Telemetry (FastAPI):} FastAPI was chosen over traditional WSGI frameworks like Flask or Django due to its native asynchronous support via ASGI (Asynchronous Server Gateway Interface). During high-intensity space shooter gameplay, the client generates telemetry JSON payloads several times per second. FastAPI handles these concurrent I/O-bound requests rapidly without blocking the Python Global Interpreter Lock (GIL), preventing upstream backpressure that could inadvertently pause or throttle the game loop.
\paragraph{} \textbf{Frontend Analytics (React):} React’s virtual Document Object Model (DOM) is optimally suited for rendering high-frequency timeseries charts. Since the \ac{dda} system modulates parameters continuously, the dashboard must redraw multi-line graphs representing enemy speed, spawn delay, and accuracy in near real-time. React limits expensive DOM reflows, ensuring the analytics layer remains highly performant even after hours of continuous gameplay logging.

\subsection{Data Pipeline Complexity}
\paragraph{} The bridge between the client-side Pygame application and the server-side FastAPI analytics layer represents a complex data engineering pipeline. The system enforces strict schema validation using Pydantic models to guarantee that malformed telemetry payloads do not corrupt the historical performance dataset. 
\paragraph{} The telemetry pipeline operates under stringent latency constraints. Specifically, the metrics collection sub-routine in the game engine must aggregate player positional deltas, calculate running health averages, and compute continuous hit-rate accuracy ratios in under 2 milliseconds. This aggregate dataset is then partitioned: one localized normalized vector is handed synchronously to the \ac{ppo} agent for immediate difficulty inference, while the raw, un-normalized dataset is serialized to a JSON payload and handed off to a daemon thread. This asynchronous daemon thread manages the HTTP POST requests to the FastAPI backend, entirely decoupling network latency from the primary game rendering loop and ensuring rigid adherence to the minimum 30 FPS non-functional requirement.

\subsection{Telemetry Data Model and API Specification}
\paragraph{} To ensure structural integrity between the diverse technology stacks, a unified JSON schema was established for all telemetry communications. The FastAPI server exposes several \ac{rest}ful endpoints, but the primary ingestion channel operates over \texttt{/api/telemetry/live} using a POST transmission. Below is the rigidly enforced data dictionary for the primary payload:

\begin{table}[!ht]
\centering
\caption{Telemetry Payload Data Dictionary}
\label{tab:data_dictionary}
\renewcommand{\arraystretch}{1.5}
\begin{tabular}{|p{3cm}|p{2cm}|p{8cm}|}
\hline
\textbf{Field Name} & \textbf{Data Type} & \textbf{Description \& Validation Bounds}\\
\hline
\texttt{session\_id} & String (UUID) & Unique identifier generated at the start of each game run.\\
\hline
\texttt{timestamp} & Float & High-precision UNIX epoch timestamp of the exact frame execution.\\
\hline
\texttt{health\_pct} & Float & Player's current health remaining. Strictly bounded $[0.0, 1.0]$.\\
\hline
\texttt{accuracy} & Float & Running ratio of successful hits versus fired lasers. Bounded $[0.0, 1.0]$.\\
\hline
\texttt{kill\_rate} & Float & Entities destroyed rolling average per second. Unbounded, typical $[0, 15]$.\\
\hline
\texttt{damage\_rate} & Float & HP lost rolling average per second. Unbounded, typical $[0, 50]$.\\
\hline
\texttt{enemy\_speed} & Float & Current dynamic scaling multiplier for enemy velocity.\\
\hline
\texttt{spawn\_delay} & Float & Current dynamic delay (in ms) between wave instantiations.\\
\hline
\texttt{ai\_action} & Integer & The raw discrete integer ($0-4$) selected by the \ac{ppo} agent.\\
\hline
\end{tabular}
\end{table}

\paragraph{} The standard API workflow dictates that the Pygame client will buffer state events until the $2.0$ second inference window triggers. At this precise moment, the dictionary is serialized, and Python's \texttt{urllib} module performs a non-blocking asynchronous POST request. If the FastAPI backend returns an HTTP 200 OK status, the local buffer is flushed. If an HTTP 500 or network timeout occurs, the client temporarily caches the payload and aggressively retries during the subsequent cycle, ensuring zero data loss during localized internet disruptions.

\section{Feasibility Study}

\subsection{Technical Feasibility}
\paragraph{} The project is technically feasible using a number of well-established technologies: Pygame for game development, PyTorch for \ac{rl} model training and inference, FastAPI for the backend \ac{api}, and React for the frontend dashboard. \ac{ppo} has been shown to be effective for continuous control and decision-making tasks in many studies. The Gymnasium library offers a standardized interface for designing \ac{rl} environments.

\subsection{Operational Feasibility}
\paragraph{} The game controls are simple and easy to learn (movement using WASD/Arrow keys, Spacebar to shoot), which are familiar to many gamers. The \ac{dda} system works transparently in the background, needing no user intervention. The analytics dashboard offers a simple web-based interface to view gameplay statistics.

\subsection{Economic Feasibility}
\paragraph{} The tools and libraries used in the project are all open-source. Python, Pygame, PyTorch, FastAPI, and React are all freely available. No proprietary datasets or paid services need to be used. As far as hardware requirements go, the system can run on standard consumer hardware with an optional GPU for faster model training.

\section{System Specification}

\subsection{Hardware Specification}
\begin{table}[!ht]
\centering
\caption{Hardware Specifications}
\label{tab:hardware}
\begin{tabular}{|p{4cm}|p{8cm}|}
\hline
\textbf{Component} & \textbf{Specification}\\
\hline
Processor & Intel Core i5 10th Generation or equivalent\\
\hline
RAM & 8 GB DDR4 (minimum)\\
\hline
GPU & NVIDIA GTX 1050 or higher (for training); CPU sufficient for inference\\
\hline
Storage & 10 GB free disk space\\
\hline
Display & 1920$\times$1080 resolution (minimum 1366$\times$768)\\
\hline
\end{tabular}
\end{table}

\subsection{Software Specification}
\begin{table}[!ht]
\centering
\caption{Software Specifications}
\label{tab:software}
\begin{tabular}{|p{4cm}|p{8cm}|}
\hline
\textbf{Software} & \textbf{Version/Details}\\
\hline
Operating System & Windows 10/11, Ubuntu 20.04+, or macOS 12+\\
\hline
Programming Language & Python 3.8+\\
\hline
Game Engine & Pygame 2.5\\
\hline
Deep Learning & PyTorch 2.0\\
\hline
RL Library & Gymnasium 0.28\\
\hline
Backend Framework & FastAPI 0.100\\
\hline
Frontend Framework & React 18\\
\hline
Data Visualization & Matplotlib 3.7, Recharts (React)\\
\hline
IDE & VS Code, Google Colab (training)\\
\hline
Version Control & Git, GitHub\\
\hline
\end{tabular}
\end{table}
